\section{Tire Physics: Relaxation \& Grip}

This section documents tire force relaxation (lag effects) and grip saturation as tires approach their limits.

\subsection{Tire Relaxation Time Constant}

\begin{equation}
  \tau_{\text{relax}}(v) = \frac{\lambda}{max(|v_{\text{long}}|, \epsilon)}
  \label{eq:relaxation-tau}
\end{equation}

\textbf{Physical Meaning}: Tires build force over a fixed distance (relaxation length $\lambda$), not a fixed time. The time constant is thus speed-dependent: $\tau(v) = \lambda / v$.

\textbf{Source Code}: \coderef{physics.js}{744--748}

\begin{lstlisting}[caption={Speed-dependent tire relaxation}]
const relaxationLength = params.tireRelaxationLength || 0.3; // metres
const vLongAbs = Math.max(Math.abs(wheelLongitudinalSpeed), 0.5);
const tireRelaxTau = relaxationLength / vLongAbs;
const tireRelaxAlpha = 1.0 - Math.exp(-dt / tireRelaxTau);
\end{lstlisting}

\textbf{Parameters}:
\begin{itemize}
  \item $\lambda = 0.3$ m: Relaxation length (typical road tire value)
  \item $\epsilon = 0.5$ m/s: Minimum speed to prevent infinite time constant
\end{itemize}

\subsection{Force Relaxation (EMA Filtering)}

\begin{equation}
  F_{\text{relaxed}} = F_{\text{prev}} + (F_{\text{current}} - F_{\text{prev}}) \cdot (1 - e^{-dt/\tau_{\text{relax}}})
  \label{eq:force-relaxation}
\end{equation}

\textbf{Physical Meaning}: Tire forces are smoothed over the relaxation time constant via exponential moving average (EMA). This creates realistic lag in the tire's response to slip changes.

\textbf{Parameters}:
\begin{itemize}
  \item $\alpha = 1 - e^{-dt/\tau}$: EMA blending factor
  \item $dt$: Physics timestep (typically 0.01 s)
\end{itemize}

\subsection{Friction Ellipse (Traction Circle)}

\begin{equation}
  \sqrt{F_{\text{lateral}}^2 + F_{\text{longitudinal}}^2} \leq N \cdot \mu
  \label{eq:friction-ellipse}
\end{equation}

\begin{equation}
  F_{\text{scaled}} = F \cdot \min\left(1, \frac{N \cdot \mu}{\sqrt{F_{\text{lateral}}^2 + F_{\text{longitudinal}}^2}}\right)
  \label{eq:friction-scaling}
\end{equation}

\textbf{Physical Meaning}: Combined lateral and longitudinal forces cannot exceed the tire's total grip circle. If they do, both forces are scaled down proportionally, preventing impossible force combinations.

\textbf{Source Code}: \coderef{physics.js}{725--736}

\begin{lstlisting}[caption={Friction ellipse clipping}]
const frictionLimit = normalLoad * frictionCoeff;
const combinedMag = Math.hypot(fadedLateralForceMag, longitudinalForceMag);
let frictionScale = 1.0;
if (combinedMag > frictionLimit && combinedMag > 0) {
  frictionScale = frictionLimit / combinedMag;
}
\end{lstlisting}

\subsection{Lateral Saturation \& Grip}

\begin{equation}
  \text{saturation} = \min\left(1, \frac{|F_{\text{lateral}}|}{N \cdot \mu}\right)
  \label{eq:lateral-saturation}
\end{equation}

\begin{equation}
  \text{grip} = \max(0, 1 - \text{saturation})
  \label{eq:grip-remaining}
\end{equation}

\textbf{Physical Meaning}: Saturation measures how much of the tire's lateral grip budget is consumed. When saturation = 1, all lateral grip is used and the tire is sliding. Remaining grip indicates steering authority available.

\textbf{Notes}: Used for visual effects (skid mark intensity) and for detecting drift state.

\newpage
